% 1- Explicar detenidamente la principal motivacion del articulo. Dejar claro cual es el posible problema
% 2- Presentar las preguntas de investigacion pertinentes
% 3- Presentar algunos resultados preliminares
% 4- Presentar el trabajo llevado hasta ahora de investigacion y desarrollo
% 5- Dejar algunas preguntas abiertas y problemas encontrados hasta ahora
%
% $Id: slides.tex 4228 2006-06-21 21:55:12Z jjamor $
%
%
% Compilar a .pdf con LaTeX (pdflatex)
% Es necesario instalar Beamer (paquete latex-beamer en Debian)
%

%
% Grficos:
% Los grficos pueden suministrarse en PNG, JPG, TIF, PDF, MPS
% Los EPS deben convertirse a PDF (usar epstopdf)
%


\documentclass{beamer}
\usetheme{Warsaw}
\usebackgroundtemplate{\includegraphics[width=\paperwidth]{format/libresoft-bg.png}}
\usepackage[english]{babel}
\usepackage[latin1]{inputenc}
\usepackage{graphics}
\usepackage{amssymb} % Simbolos matematicos
\usepackage{url}

%\definecolor{libresoftgreen}{RGB}{162,190,43}
%\definecolor{libresoftblue}{RGB}{0,98,143}

%\setbeamercolor{titlelike}{bg=libresoftgreen}

%% Metadatos del PDF.
\hypersetup{
  pdftitle={Libresoft},
  pdfauthor={Daniel Izquierdo Cortazar},
  pdfcreator={GSyC/Libresoft},
  pdfproducer=PDFLaTeX,
  pdfsubject={Master on Free Software},
}
%%



\AtBeginSection[]
{
  \begin{frame}<presentation>
    \frametitle{Index}
    \tableofcontents[current]
  \end{frame}
}




\begin{document}



\title{Introduction to \textit{Evaluation of Libre Software Projects}
}
\subtitle{Master on Free Software
}
\author{Felipe Ortega, Pedro Coca, Daniel Izquierdo}
\date{Universidad Rey Juan Carlos, Fuenlabrada, Spain\\ 18th, November, 2010}

\frame{
\maketitle
\begin{center}
\includegraphics[width=4cm]{format/gsyc-urjc.pdf} 
\end{center}
}


% Si el titulo o el autor se quieren acortar para los pies de pgina
% se pueden redefinir aqu:
\title{Evaluation of Libre Software Projects}
%\author{Daniel Izquierdo}


%% LICENCIA DE REDISTRIBUCION DE LAS TRANSPAS
\frame{
\vspace{4cm}

\begin{flushright}
{\tiny
(cc) 2010 Daniel Izquierdo Cort\'azar. \\
Some rights reserved. This document is distributed under the Creative \\
            Commons Attribution-ShareAlike 2.5 licence, available in \\
            http://creativecommons.org/licenses/by-sa/2.5/

%  Este documento (o uno muy similar) está disponible en \\
%  \url{http://gsyc.escet.urjc.es/~jjamor/}
}
\end{flushright}
}
%%

\frame{
\frametitle{Table of contents}
\tableofcontents
}


\section{Admiministrative Data}

\begin{frame}
\frametitle{Administrative data}

\begin{itemize}
\item First semester, 3 ECTS credits
\item Teachers:
  \begin{itemize}
  \item Felipe Ortega (jfelipe at gsyc.es)
  \item Pedro Coca (pcoca at gsyc.es)
  \item Daniel Izquierdo (dizquierdo at gsyc.es)
  \item Departamento de Sistemas Telemticos y Computacin (GSyC)
  \item Rooms 106, 118 and 120 Departamental II (Mostoles campus)
  \end{itemize}
\end{itemize}
\end{frame}



\begin{frame}
\frametitle{Administrative data}

\begin{itemize}
\item Schedule: Thursday, 17:00-21:00
\item Sessions:
  \begin{itemize}
  \item Classroom 215, Aulario II, Fuenlabrada campus
  \item Lab 7, Laboratorios II, Fuenlabrada campus
  \end{itemize}
\item Moodle course (please, join it as soon as possible): \\
  \url{http://gsyc.escet.urjc.es/moodle/course/view.php?id=125}
\end{itemize}
\end{frame}

\section{Subjects}

\begin{frame}
\frametitle{Main topics addressed}

\begin{itemize}
\item Introduction to Quality from a Software Engineering perspective
\item Quality models in Open Source: OpenBRR, QSoS
\item Quality models in Open Source: QualOSS, SQO-OSS
\item Automation of quality models
\end{itemize}

\begin{quote}
\begin{center}
Provide a framework for analyzing and evaluating libre software projects
\end{center}
\end{quote}

\end{frame}

\section{Evaluation}


\begin{frame}
\frametitle{Evaluation}

\begin{itemize}
\item Exercises
\item Blog entries
\item Collaborative work
\item Specific report
\item ...
\end{itemize}

\end{frame}



\section{References}

\begin{frame}
\frametitle{Some references}

\begin{itemize}
\item Sitio web de la asignatura \\
  \url{http://docencia.etsit.urjc.es/moodle/course/view.php?id=125}

\end{itemize}

\end{frame}


\end{document}
